\section{SOS Save Our Souls}

\begin{quotex}
Every unhappy family is unhappy in its own way. 
\flright{\emph{Anna Karenina}}
\end{quotex}


The international distress signal is dot dot dot dash dash dash dot dot dot, or SOS for short. It is easy to remember as ``Save our ship'' or ``Save our souls''. In the latter case, there should be many more such signals.

At least SOS is organic, and Morse code adepts could identify the sender based on his style. 911, as the emergency number, is not so; the best I could come up with is ZAB, but it may never catch on. It is not even international. Europeans have BAD instead.

Apart from the rare cases of those predestined to righteousness, most people seek out a spiritual path following a boundary situation or an inner collapse from suffering, guilt, illness, anxiety, despair, and the like. The majority of them will gravitate to an exoteric religion or its equivalent. They want an easy way to save their souls from that collapse. Moreover, they are attracted to promises of a pretty wife or rich husband, a nice car, a better job, perfect health, etc.

Those who are naturally blessed with such things are seldom so motivated. Or else they treat their spirituality as if it were a pretty ornament, to add some ``style'' to their characters. Their spirituality is ``bought and paid for''; they pay for a shaman, a Jungian analyst, a weekend seminar, a course on the ``law of attraction'', and so on. Unfortunately, what they want cannot be bought.

\paragraph{Escapism}


Those who choose a path are often accused of ``escapism''. Clearly, when the house is on fire, is not escape a sound instinct? Real escapism is pretending the danger does not exist. Such people think that the fire is still far away, or there will always be time later on.

The Buddha famously came to the realization that all life is suffering. He therefore gave up his family, his throne, his wealth, his power in order to escape from that suffering. What did he see that few in his position can see? Nothing significant has changed.

\paragraph{Moment of Decision}


An esoteric path, unlike an exoteric path, promises something different. There is an effort required to ``save one's soul''. This means that the crucifixion and descent into hell comes before resurrection. For them, salvation and liberation is achieved by developing a Real Self, Consciousness, and a True Will. That has little appeal because people assume they already have those things.

Hermetists are indistinguishable\footnote{\url{https://gornahoor.net/?p=793}} by sight. That is why they used to pose as horse traders or street performers, etc. Anything to give them cover, and freedom of movement, for their real task. They will appear, seemingly randomly, at various times and places. Those who truly desire an esoteric path will seize on those opportunities, since they are rare and hard to find. Yet most men are stragglers. They dream of such a path, yet refuse to actually enter on the Way. They have a consumer mentality and believe that such paths will always be available to them, since the supermarkets have an abundance of milk and iPhones are everywhere.

\begin{itemize}


\item Cult: easy to find, hard to leave.

\item Esoteric way: hard to find, easy to leave
\end{itemize}


\paragraph{Signs of the Times}


There is still time left until we hit bottom. As of 2014 AD we are in the seventh circle of hell. The signs can be recognized from the following stages. What was forbidden gets tolerated. Then it becomes normalized. Finally, it is considered praiseworthy. Look it up.

Suicide is gaining respectability. Those in favor of it desire access to an effective and pain free biochemical death. Yet those who oppose capital punishment claim such a death is not possible, so it is cruel and unusual punishment. Based on that alone, you might suppose that the Venn diagram for pro-suicide and anti-capital punishment activists would not intersect. However, I suspect that they very nearly coincide.

Blasphemers and atheists sell a lot of books.

Usurers make the big money. For them, money is no longer a medium of exchange of products, but the product itself. It would be like selling an ``inch'', the inch of what does not matter.

The latest craze is the legalization of marijuana, particularly among the best and the brightest. I suppose that 50,000 years of evolution have blessed them with such large brains, that they now exercise the option not to use them.

Many women claim that their dog, or dogs, is her best friend. She is absolutely convinced that the dog understands her. Yet she is alone and single because her husband or boyfriend evidently did not understand her.

Anti-Christian is code word for anti-European, since it ignores and denies 1500 years of High Culture. Those who would destroy a people, do so by effacing their history.

Christianity is esoteric paganism. It does not destroy previous traditions, but incorporates them into a larger whole. Yet paganism survives, e.g., in images of God as a being rather than Being. This image is like Zeus and Hera looking over the human race and arbitrarily interfering in events from time to time, picking winners and losers. That is idolatry. You can remove statues and icons from so-called churches, but as long as there is an image of God in your mind, that is still idolatry from the esoteric perspective.

\begin{quotex}
Happy families are all alike.
\end{quotex}



\flrightit{Posted on 2014-11-07 by Cologero}

\begin{center}* * *\end{center}

\begin{footnotesize}
\begin{sffamily}
\texttt{Wayne Ferguson on 2014-11-07 at 08:43 said:} 

``Suicide is gaining respectability. Those in favor of it desire access to an effective and pain free biochemical death. Yet those who oppose capital punishment claim such a death is not possible, so it is cruel and unusual punishment. Based on that alone, you might suppose that the Venn diagram for pro-suicide and anti-capital punishment activists would not intersect. However, I suspect that they very nearly coincide.'' But are there not important distinctions to be drawn between Brittany Maynard (the young woman in the news this past week) and an inmate on death row? Is it really so strange that the same physical procedure might be judged humane in one case but not in the other?

\hfill

\texttt{Cologero on 2014-11-07 at 08:49 said:}

Wayne, I was not opening up a debate about capital punishment and suicide. That paragraph is intended to be understood in the context of Dante's seventh circle of hell. BTW, I do think it is really strange.

\hfill

\texttt{Cologero on 2014-11-07 at 09:21 said:}

Just to provide some context that I left out but had in the back of my mind. Several months ago a biochemical execution in Oklahoma went horribly wrong. That was hardly humane. Bracketing out the question of the man's guilt and the inhumanity of his crime, the US constitution prohibits cruel punishments. I doubt euthanasia supporters would want the same sort of self-execution.

\hfill

\texttt{Wayne Ferguson on 2014-11-07 at 09:23 said:}

I'm not interested in debating these topics either--just pointing out that those critical of capital punishment and sympathetic to voluntary euthanasia may consistently hold both positions (the physical and pychological state of the person in question being decisive to the question of the humaneness of the procedure as they see it).

With regard to the seventh circle of hell, it seems that the suicide deserves to suffer more than our contemporary ``violent offenders'' (assuming that the occupants of the middle circle are guilty of a more profound offense than those in the outer circle). Is that right?

\url{http://en.wikipedia.org/wiki/Inferno\_(Dante)\#Seventh\_Circle\_.28Violence.29}


\hfill

\texttt{scardanelli on 2014-11-07 at 10:26 said:}

``Christianity is esoteric paganism. It does not destroy previous traditions, but incorporates them into a larger whole.''

For those who may still doubt this or better, have yet to experience it, David Fideler's Jesus Christ, Sun of God is an excellent resource. I'm only halfway through it, but it has thus far quite solidly demonstrated the continuation of ancient tradition in early Christianity... in the conception of the Logos, in the parables of the feeding of the 5000 and the 153 fishes in the unbroken net, etc. Greek geometry, gematria, and the Pythagorean tradition are all present in the gospels, leading one to believe that the authors of the Gospels were well aware of it.

\hfill

\texttt{Wayne Ferguson on 2014-11-07 at 11:09 said:}

Thanks for mentioning David Fideler's book--found a copy at Scribd:

\url{https://www.scribd.com/doc/230403079/Jesus-Christ-Sun-of-God-Ancient-Cosmology-and-Early-Christian-Symbolism-by-David-R-Fideler}


\hfill

\texttt{X on 2014-11-07 at 19:16 said:}

SYOSF Save Your Own Soul First

\hfill

\texttt{White Rabbit on 2014-11-08 at 02:47 said:}

I can see why many women believe dogs understand them. Today I met a retarded person. I was doing some public sloganeering and noticed he was responsive. I looked at him, and upon finding his speech incomprehensible and seeing a giant grin on his face realized that he was mentally handicapped. As I shook his hand and briefly lectured to him while he stood there with his beaming smile, I still felt that he was absorbing what I was saying as. Blank walls are fit for projection.

\hfill

\texttt{August on 2014-11-09 at 06:26 said:}

That's a good point White Rabbit.

A close reading of posts here will reveal that most people are passive, subject to directed influences. If your sloganeering meets with little success, we can safely assume that those you cross paths with are already subordinated to some other more formidable impression.

Though unlike your retard they have a special blessing/curse, in that they can be convinced that they are entirely free and autonomous actors.

\hfill

\texttt{argusandphoenix on 2014-11-17 at 10:37 said:}

This is sociology, but it's interesting and apt as an example of programming:

\url{http://fabiusmaximus.com/2014/11/16/castle-tv-plot-women-rights-feminism-72870/}


\end{sffamily}
\end{footnotesize}

